\section{From supersaturation to a finite graph}\label{sec:distributions}

A \emph{unit} is an ordered pair of raw vertices with no hole between its endpoints.
The endpoints of a unit may be dependent.
When two units are sampled independently, their four endpoints need not be independent within either unit.
The event of interest is that all four cross pairs have holes, as in \Cref{fig:four-holes}.

% lax begin Lax793645.FourHoles
\begin{theorem}[Raw supersaturation]\label{thm:raw-supersaturation}
For the construction parameters specified in \Cref{sec:geometry,sec:parameters}, for all sufficiently large \(n\), every probability law \(\sigma\) on units satisfying
\begin{equation}\label{eq:raw-law-caps}
 \sigma_1\le M\mu,\qquad \sigma_2\le M\mu,\qquad
 \sigma\le 2^{DN}\mu^2
\end{equation}
gives probability at least \(2^{-200N}\) that two independent units have all four cross holes.
The inequalities in \Cref{eq:raw-law-caps} are pointwise inequalities of measures on the finite raw spaces.
\end{theorem}
% lax end

The proof of \Cref{thm:raw-supersaturation} occupies the subsequent sections and is completed in \Cref{sec:rare-overlap}.
We prove its implication for graphs here. The frame estimates and the preparation of an admissible unit law begin in \Cref{sec:frame-laws}.

\subsection{A finite exposure potential}

Join two raw units when all four cross pairs are holes. The resulting
conflict graph has no loops. Write $q=\mu^2$,
$\kappa=2^{DN}$ and $\varepsilon=2^{-200N}$.
For $\gamma>0$ and an integer $b>0$ with $\kappa\le2^b\gamma$, put
\[
 \Phi(\sigma)=\sum_{i=0}^{b-1}\sum_x
      \min\{\sigma(x),2^i\gamma q(x)\}.
\]
Every summand indexed by $i$ has total at most one, so
$0\le\Phi(\sigma)\le b$ for every probability law $\sigma$.

% lax begin Lax793645.ClippedExposure
\begin{lemma}[Clipped exposure]\label{lem:clipped-exposure}
Let $(Y,X)$ have both marginals $\sigma$ and never consist of
conflicting units. Write $\tau_y$ for the conditional law of $X$
given $Y=y$. If $\sigma\le\kappa q$ on a set $A$, then
\[
 \Phi(\sigma)-\E_Y\Phi(\tau_Y)
 \ge \frac{(\sigma\otimes\sigma)(\text{conflict}\cap A^2)-\gamma}{2}.
\]
The potential drop is always nonnegative, without any assumption on $A$.
\end{lemma}
% lax end
% lax begin Lax793645Proofs.Manuscript.ClippedExposure.dyadic_exposure_joint
\begin{proof}
For each $x$ and each threshold $a=2^i\gamma$, the identity
$\E_Y\tau_Y(x)=\sigma(x)$ gives
\[
 \E_Y\min\{\tau_Y(x),aq(x)\}
 \le\min\{\sigma(x),aq(x)\}.
\]
Indeed the left side is separately bounded by both entries of the
minimum. Let $p_x$ be the probability that $Y$ conflicts with $x$. On that event,
$\tau_Y(x)=0$. If $\gamma q(x)\le\sigma(x)\le\kappa q(x)$,
choose a threshold with $aq(x)\le\sigma(x)\le2aq(x)$. The drop
at this coordinate and threshold is at least
$p_xaq(x)\ge p_x\sigma(x)/2$. Coordinates below $\gamma q(x)$
have total mass at most $\gamma$. Sum over $x\in A$ and discard
all other nonnegative drops. Restricting the first unit to $A$ can
only decrease the conflict probability.
\end{proof}
% lax end

\subsection{Sampling repeatedly from a matching}

% lax begin Lax793645.MatchingExposure
\begin{proposition}\label{prop:raw-to-graph}
Assume \Cref{thm:raw-supersaturation} for a sufficiently large $n$.
There is a sample of $m=2^{2000N}$ raw vertices whose graph satisfies
\[
 \alpha(G)\le2,\qquad \cm(G)<m/100.
\]
\end{proposition}
% lax end
% lax begin Lax793645Proofs.Manuscript.MatchingExposure.transfer_scales
\begin{proof}
Suppose instead that every sample has a touching matching of size
$k=\lceil m/100\rceil$. Choose one such matching deterministically
for each sample, and orient its edges arbitrarily. Given the sample,
draw repeatedly and independently from these $k$ oriented edges.
Two emissions never conflict, even if they are the same edge.
The hidden state is the original sample, with law $\mu^m$.

After any history $h$, let $\sigma_h$ be the law of the next emitted
unit under the posterior law of the hidden sample. Two future
emissions have both marginals $\sigma_h$ and never conflict.
The law after observing $y$ is the conditional second marginal
required by \Cref{lem:clipped-exposure}.

Define
\[
 S_h=\{x:(\sigma_h)_1(x)>M\mu(x)/4
               \text{ or }(\sigma_h)_2(x)>M\mu(x)/4\},
 \qquad
 E_h=\{x:\sigma_h(x)>\kappa q(x)/4\}.
\]
Summing each density excess gives
$\mu(S_h)\le8/M$ and $q(E_h)\le4/\kappa$.
Let $A_h$ consist of units avoiding $S_h$ at both endpoints and
avoiding $E_h$, and set $\alpha_h=\sigma_h(A_h)$.
If $\alpha_h\ge1/4$, the normalized restriction to $A_h$ satisfies
all three caps in \Cref{eq:raw-law-caps}. Raw supersaturation then gives
\[
 (\sigma_h\otimes\sigma_h)(\text{conflict}\cap A_h^2)
 \ge\varepsilon\alpha_h^2\ge\varepsilon/16.
\]
Use $\gamma=\varepsilon/32$. The drop $D_h$ in the potential is at
least $\varepsilon/64$ on these histories and is always nonnegative.
Since $\alpha_h\le1$, in all cases
\begin{equation}\label{eq:exposure-retention}
 D_h\ge\frac{\varepsilon}{48}(\alpha_h-1/4).
\end{equation}

Take $L=2^{400N}$ exposure steps and $s=\lfloor k/5\rfloor$.
There are at most $(|\Omega_n|^2+1)^L$ histories of length at most
$L$: this follows by induction from the recursion that either stops
or appends one unit. Each $S_h,E_h$ is a deterministic set in the raw
space. The iid sampling bounds therefore apply before any exposure.
A union bound shows that the probability of a raw sample having,
for some such history, $s$ positions in $S_h$ or $s$ disjoint
position pairs in $E_h$ is at most
\begin{equation}\label{eq:exposure-exceptions}
 \Delta=(|\Omega_n|^2+1)^L
 \left[2^m(8/M)^s+m^{2s}(4/\kappa)^s\right].
\end{equation}
Call this fixed event of raw samples bad. On a sample outside it,
at most $2(s-1)$ edges of the chosen matching fail the restrictions
for any history. Since $4s\le k$, its emission law gives $A_h$
probability at least $1/2$.

Let $B_h$ be the posterior probability of the fixed bad event.
Averaging the preceding bound under the original posterior gives
$\alpha_h\ge1/2-B_h$. The posterior probability $B_h$ is a
martingale: averaging its values after the next observation returns
its current value. Thus \eqref{eq:exposure-retention} and a finite
tree telescope imply
\[
 b\ge L\left(\frac{\varepsilon}{192}
                  -\frac{\varepsilon}{48}\Pr(\text{bad})\right).
\]
No posterior law has been conditioned on the complement of the bad event.

Here $D=8000$, $M=2^{1000}$ and $m=2^{2000N}$.
We may take $b=8200N+6$, so $\kappa\le2^b\gamma$.
By \Cref{eq:raw-space-size}, for all sufficiently large $N$,
\[
 (|\Omega_n|^2+1)^L\le2^{m/2},\qquad m\le550s.
\]
The two summands in \eqref{eq:exposure-exceptions}, including this
history factor, are each at most $2^{-m/4}$. For the first, compare
$m/2+m-997s$ with $-m/4$; for the second, compare
$m/2+(2-4000N)s$ with $-m/4$. Consequently $\Delta<1/8$.
But $L\varepsilon=2^{200N}>384b$ for large $N$, contradicting
the telescope. The required sample exists. Finally,
$\alpha(G)\le2$ holds for every sample by \Cref{lem:hole-triangle-free}.
\end{proof}
% lax end

\subsection{The clique-minor bound}

\Cref{prop:raw-to-graph} supplies the required bound on connected
matchings once raw supersaturation is proved. The next counting
argument converts that bound into the clique-minor estimate used in
the introduction.

The following elementary estimate counts small branch sets, in the
same connected-matching approach used by Cambie~\citep{cambie2021}.
Its deliberately coarse bound holds for every finite graph.

% lax begin Lax793645.MatchingMinorBound
\begin{proposition}\label{prop:matching-minor}
For every finite nonempty graph \(G\) of order \(m\),
\[
 \had(G)\le\frac{m+4\cm(G)+2}{3}.
\]
In particular, if \(\alpha(G)\le2\) and \(\cm(G)<m/100\), then
\[
 \had(G)<\frac{26m}{75}+\frac23<\frac m2\le\chi(G)
 \qquad(m\ge5).
\]
\end{proposition}
% lax end
% lax begin Lax793645Proofs.HadwigerMatching.hadwiger_matching_bound
\begin{proof}
Consider a complete-minor model with \(b\) branch sets.
Let \(s\) of them be singletons and \(e\) have two vertices.
Each two-vertex branch is an edge.
The singleton vertices form a clique, so pairing them gives \(\lfloor s/2\rfloor\) more edges.
All these edges together form one touching matching: their vertex sets are disjoint, and adjacency of the original branches guarantees every required contact.
Writing \(c=\cm(G)\), we obtain
\[
 c\ge e+\lfloor s/2\rfloor,\qquad s\le2(c-e)+1.
\]
Every remaining branch has at least 3 vertices, whence
\[
 m\ge s+2e+3(b-s-e)=3b-2s-e.
\]
Combining the inequalities gives
\[
 3b\le m+2s+e\le m+4c-3e+2\le m+4c+2.
\]
Maximize over all complete-minor models.
For the final assertion, every color class has at most 2 vertices when \(\alpha(G)\le2\), so \(\chi(G)\ge m/2\).
Combine this with \(c<m/100\) and \(26m/75+2/3<m/2\) when \(m\ge5\).
\end{proof}
% lax end
